\newpage
\chapter{Mathematical Optimization}

\emph{Optimization} associates a cost function \(f:\R^n \to \R^n\) with a singular or multiple optimal 
solution \(x \in \R^n\). They can also be associated with some constraints.

Optimization relates closely to derivatives and iterative numerical methods for finding the solution \(x\).

\section{Extrema and Critical Points Test Of Single Variable Calculus}

A function \(f\) has a critical point at \(x = c\) if:

\[
    f'(c) = 0 \quad \text{or} \quad f'(c) \text{ does not exist}
\]

To determine the nature of the critical point:

\begin{itemize}

    \item \emph{First Derivative Test:} Check sign changes of \(f'\) around \(c\)

    \item \emph{Second Derivative Test:}
    \[
        f''(c) > 0 \Rightarrow \text{local minimum} \\
        f''(c) < 0 \Rightarrow \text{local maximum}
    \]

    This works, because the second derivative is telling us about the concavity of the function. If it is 
    positive then we will have function which is concave up a minimum; and if negative concave down a maximum. 

\end{itemize}

\section{Relative Extrema In \texorpdfstring{\(R^n\)}{}}

To determine local extrema for \( f : \R^n \to \R \), follow this procedure:

\textbf{Step 1:} Find critical points by solving:

\[
    \nabla f(\vec{x}) = 0
\]

\textbf{Step 2:} Compute the Hessian matrix:

\[
    H_f(\vec{x}) = \left( \frac{\partial^2 f}{\partial x_i \partial x_j} \right)
\]

\textbf{Step 3:} Analyze the Hessian at critical points:

\begin{itemize}

    \item If \( H \) is positive definite \(f_{xx}(x_0, y_0) > 0, \det(H) > 0 \implies\) local minimum

    \item If \( H \) is negative definite \(f_{xx}(x_0, y_0) < 0, \det(H) > 0 \implies\) local maximum

    \item If \( H \) has mixed signs (indefinite) \(\det(H) < 0 \implies  \) saddle point

    \item If \(\det(H) = 0 \implies \) Next derivative decides

\end{itemize}

The determinant is needed as part of the test because, while the second derive tells us about the concavity, the 
determinant help us determine if all the concavities given by the axis second derivatives with respect to 
\(x_1, x_2, \dots, x_n\) and the mix partials are up, down or neither. If the signs are mixed then it is a saddle.  

\textbf{Example: \( f(x, y) = x^2 + y^2 \)}

\[
    \nabla f = \begin{bmatrix} 2x \\ 2y \end{bmatrix} \Rightarrow \text{Critical point at } (0, 0)
\]

\[
    H_f = \begin{bmatrix} 2 & 0 \\ 0 & 2 \end{bmatrix} \Rightarrow \text{positive definite}
    \Rightarrow (0, 0) \text{ is a local minimum}
\]

\section{Finding Candidate Points for Zeros While Seeking Relative Extrema}

To find local minima or maxima of a function \( f : \R^n \to \R \), we first identify 
\emph{critical points}, which are the candidates for relative extrema. 
These are the points where the gradient of \(f\) vanishes or does not exist.

\textbf{Step 1: Compute the Gradient}

The gradient of \(f\), denoted \( \nabla f \), collects all the first partial derivatives:

\[
    \nabla f(x_1, \dots, x_n) =
    \begin{bmatrix}
    \frac{\partial f}{\partial x_1} \\
    \vdots \\
    \frac{\partial f}{\partial x_n}
    \end{bmatrix}
\]

\textbf{Step 2: Solve the System \( \nabla f = 0 \)}

Find all points \( \vec{x}_0 \in \R^n \) such that:

\[
    \nabla f(\vec{x}_0) = \vec{0}
\]

This is a nonlinear system of \(n\) equations in \(n\) variables. The solutions \( \vec{x}_0 \) 
are the \emph{candidate points for extrema}, and also referred to as the 
\textbf{stationary points} or \emph{critical points} of the function.

You use the following techniques to find the points: 

\begin{itemize}

    \item Factorization

    \item Addition of rows

    \item Substitution

\end{itemize}

There are more but which one to use depends on the kind problem.

\textbf{Step 3: Analyze the Nature of the Critical Points}

Once all critical points \( \vec{x}_0 \) satisfying \( \nabla f(\vec{x}_0) = 0 \) are found, analyze 
each using the \emph{second derivative test} involving the Hessian matrix \( H_f \):

\[
    H_f(\vec{x}) = \left( \frac{\partial^2 f}{\partial x_i \partial x_j} \right)
\]

Evaluate \( H_f \) at each critical point and use the eigenvalues or definiteness to classify:

\begin{itemize}

    \item If \( f_{xx}(x_0, y_0) > 0 \) and \( D > 0 \), then \(f\) has a \emph{local minimum} at 
          \( (x_0, y_0) \).

    \item If \( f_{xx}(x_0, y_0) < 0 \) and \( D > 0 \), then \(f\) has a \emph{local maximum} at 
          \( (x_0, y_0) \).

    \item If \( D < 0 \), then \( (x_0, y_0) \) is a \emph{saddle point}.

    \item If \( D = 0 \), the \emph{second derivative test is inconclusive}; higher-order derivatives 
          must be considered.

\end{itemize}

\subsection{Non-differentiable Points}

If the function \(f\) is not differentiable at some point \( \vec{x}_0 \), but defined there, this 
point may also be a candidate for extremum and should be examined separately using limit-based analysis 
or directional behavior.

\section{Lagrange Multipliers}

Let \( f : \R^n \to \R \) be the function to optimize, subject to the constraint 
\( g(\vec{x}) = 0 \), where \( g : \R^n \to \R \).  

If we look at the constraint intersection with the original function we are going to notice that along 
the intersection there are maxima and minima with their respective tangent lines and gradients that are 
normal to them. We also are going to see that both the gradient of \(f\) and the gradient of \(g\) are just scalar version of each other so: 

This leads to the \emph{Lagrange system}:

\[
    \begin{cases}
        \nabla f(\vec{x}) = \lambda \nabla g(\vec{x}) \\
        g(\vec{x}) = 0
    \end{cases}
\]

For more conditions  we can use \(g_1, \dots, g_n\) and \(\lambda_1, \dots, \lambda_n\)

\[
    \begin{cases}
        \nabla f(\vec{x}) = \lambda_1 \nabla g_1(\vec{x}) + \cdots + \lambda_n \nabla g_n(\vec{c}) \\
        g_1(\vec{x}) = 0 \\
        \vdots \\
        g_n(\vec{x}) = 0
    \end{cases}
\]

We can also compact these equations in the \emph{Lagrangian}

\[
    \mathcal{L}(x,y, \dots, \lambda_1, \dots, \lambda_n) = f(x, y, \dots) + \lambda_1 g_1(x, y, \dots) + 
    \cdots + \lambda_n g_n(x, y, \dots)
\]

\subsection{Determinant method}

Sometimes it is a good way to solve a problem of this topic is to:

\begin{enumerate}
    
    \item Build the \emph{Lagrangian}
    
    \item \(\nabla L = 0\) and maybe compute the determinant to generate another equation to solve
        
    \item Get rid of \(\lambda\) via factorization, substitution, addition of rows, etc.
    
    \item Solve the system without \(\lambda\)

\end{enumerate}

\section{The Bordered Hessian}

The Bordered Hessian matrix is:

\[
    H_B =
    \begin{bmatrix}
    0 & {\left( \nabla g(x_1, \dots, x_n) \right)}^T \\
    \nabla g(x_1, \dots, x_n) & \nabla^2 \mathcal{L}(x, \dots, x_n, \lambda_1, \dots, \lambda_n)
    \end{bmatrix}
\]

Where:

\begin{itemize}

    \item \( \nabla  g(x_1, \dots, x_n) \) is the \( m \times n \) Jacobian matrix of the constraints.

    \item \( \nabla^2 \mathcal{L}(x, \dots, x_n, \lambda_1, \dots, \lambda_n) \) is the \( n \times n \) 
          Hessian of the Lagrangian with respect to \(x\).

\end{itemize}

\textbf{Example I:}

Find the point on the circle \(x^2 + y^2 = 4\) that is the closest to the point \((3,4)\).

The function \(f(x,y)\) is this case the euclidean norm:

\[
    f(x, y) = \sqrt{(x - 3)^2 + (y - 4)^2}
\]

subject to the constraint:

\[
    g(x, y) = x^2 + y^2 - 4 = 0
\]

The gradient of the constraint is:

\[
    \nabla g = \begin{bmatrix} 2x \\ 2y \end{bmatrix}
\]

Since \( \nabla g \ne 0 \), the rank condition is fulfilled, and no further critical points 
need to be analyzed.

\textbf{Step 1: Build the Lagrangian}

\[
    \mathcal{L}(x, y, \lambda) = f(x, y) + \lambda g(x, y)
\]

In this case we can use a property of the root the function to our advantage. Notice that the minimum 
of the root function is the minimum of the term under root therefore, we can simplify some work by writing 
\(f(x,y)\) without the root. 

\[
    \mathcal{L}(x, y, \lambda) = (x - 3)^2 + (y - 4)^2 + \lambda (x^2 + y^2 - 4)
\]

\textbf{Step 2: Compute the necessary conditions}

\begin{align*}
   \textbf{I } &\frac{\partial \mathcal{L}}{\partial x} = 2(x - 3) + 2\lambda x = 0 \\
   \textbf{II } &\frac{\partial \mathcal{L}}{\partial y} = 2(y - 4) + 2\lambda y = 0 \\
   \textbf{III } &\frac{\partial \mathcal{L}}{\partial \lambda} = x^2 + y^2 - 4 = 0
\end{align*}

Now, this is the difficult part. The best approach in this case is to use substitution to 
with a top-bottom approach:

\[
    x = \frac{3}{1 + \lambda}
\]

\[
    y = \frac{4}{1 + \lambda}
\]

Solving for \(\lambda\) is also an option but in these kinds of systems of equations where 
both the partial with respect to \(x\) and \(y\) are very similar it is a good idea to solve for them 
because they tend to have a similar form when solving for them.

\begin{align*}
    \left(\frac{3}{1 + \lambda}\right)^2 + \left(\frac{4}{1 + \lambda}\right)^2 - 4 &= 0\\
    \frac{9}{(1 + \lambda)^2} + \frac{16}{(1 + \lambda)^2} - 4 &= 0\\
    \frac{25}{(1 + \lambda)^2} &= 4 \\
    25 &= 4 (1 + \lambda)^2 \\
    \sqrt{\frac{25}{4}} &= 1 + \lambda \\
    \lambda &= -1 \pm \frac{5}{2}
\end{align*}

\[
    \lambda_1 = \frac{-7}{2} \quad \lambda_2 = \frac{3}{2}
\]

Now, let us substitute backwards for \(\lambda_1\):

\[
    x = \frac{3}{1 - \frac{7}{2}} = \frac{6}{-5}
\]

\[
    y = \frac{4}{1 - \frac{7}{2}} = \frac{8}{-5}
\]

Now,  for \(\lambda_2\):

\[
    x = \frac{3}{1 + \frac{3}{2}} = \frac{6}{5}
\]

\[
    y = \frac{4}{1 + \frac{3}{2}} = \frac{8}{5}
\]

\textbf{Step 3: Hessian Test with Constraint}

We construct the Hessian:

\[
    H = \begin{bmatrix}
    2 + 2\lambda & 0 & 2x \\
    0 & 2 + 2\lambda & 2 \\
    2x & 2 & 0
    \end{bmatrix}
\]

For \(\lambda_1\) we get \( \det(H) =  \frac{244}{5} > 0 \), a minimum 
and for \(\lambda_2\) \(\det(H) = \frac{-244}{5}\) a maximum.

\textbf{Example II:}

Optimize the distance to the origin \(f(x,y,z)\) with respect to the constraints 
\(g_1(x,y,z) = z^2 - x^2 - y^2 = 0\) and \(g_2(x,y,z) = x - 2z - 3 = 0\)

\begin{align*}
    \nabla f &= \lambda \nabla g_1 + \mu \nabla g_2 \\
    g_1 &= 0 \\
    g_2 &= 0 
\end{align*}

We know that \(f(x,y,z) = \sqrt{x^2 + y^2 + z^2}\). Like with other minimization problems we can 
just find the minimum of the part under root thus, \(f(x,y,z) = x^2 + y^2 + z^2\). Now we can put 
the build our system of equations

\begin{align*}
    \textbf{I   } &2x = \lambda(-2x) + \mu \\
    \textbf{II  } &2y = \lambda(-2y) \\
    \textbf{III } &2z = \lambda(2z) + \mu(-2) \\
    \textbf{IV  } &z^2 = x^2 + y^2 \\
    \textbf{V   } &x = 2z + 3   
\end{align*}

For the equation \textbf{II} we can already see that the equation is equal to zero only if 
\(\lambda = -1\) or \(y = 0\). Thus, we have got our first numbers to work with.

\textbf{Case \(\lambda = -1\):}

\begin{align*}
    \textbf{I   } &2x = -1(-2x) + \mu  \implies \mu = 0\\
    \textbf{III } &2z = -1(2z) + 0(-2) \implies z = 0 \\
    \textbf{IV  } &0^2 = x^2 + y^2  \implies x = y = 0\\
    \textbf{V   } &0 = 0 + 3 \implies 0 = 3   
\end{align*}

With the last contradiction we know that the whole case \(\lambda = -1\) is not valid.

\textbf{Case \(y = 0\):}

\begin{align*}
    \textbf{IV  } &z^2 = x^2 + y^2 \implies z^2 = x^2  \implies z = \pm x\\
    \textbf{V   } &x = 2x + 3  \implies x = -3 (z = x)\\ 
    \textbf{V   } &x = -2x + 3  \implies x = 1 (z = -x)\\ 
\end{align*}

Now with this new information we can get two different points \(P_1(1, 0, -1)\) and \(P_2(-3, 0, -3)\).
We do not need to write the Hessian in this case because with geometric intuition the problem can be 
solved.

\[
    f(P_1) = 1^2 + 0^2 + (-1)^2 = 2 
\]

So the distance is \(\sqrt{2}\).

\[
    f(P_2) = -3^2 + 0^2 -3^2 = 18
\]

Thus, the distance is \(3\sqrt{2}\).

Our final Answer is that \(P_1\) is the closest point to the intersection of the two planes with the 
function \(f\).

\section{Convexity}

A function \(f\) is \emph{convex} is \(\metric{x - y} \in X, \ \forall x, y \in X\) where \(X\) is the \emph{epigraph} of the function. Visually 
in 2D and 3D as cup shaped. 

\emph{Convexity} delivers the property that local extrema become global extrema. In the context of optimization \emph{convexity} generalizes 
linearity.

\section{Penalty Functions}

A way to get rid of constraints in an optimization problem is to add the constraint equality or inequality as a part of our target 
function \(f\) accompanied by a scalar for the linear case and more exotic terms for non-linear cases. This technique allows for the maintenance 
of continuity.

\section{The Primal And Dual Problem}

Given a convex optimization problem with a target function \(f(X)\), equality constraints \(h_i(X)\), and inequality constraints \(g_i(X)\) we 
define the \emph{primal problem} as 

\[
    \min_{X \in \R^n} f(X) + \max_{v_i, u_i \ge 0} \sum_i v_i h_i(X) + \sum_j u_j g_j(X)
\]

with \(u_i > 0\) and \(v_i \ne 0\). We can build as the \emph{dual problem} as

\[
     \max_{v_i, u_i \ge 0}  \min_{X \in \R^n} f(X) + \sum_i v_i h_i(X) + \sum_j u_j g_j(X)
\]

this problem can be easier to solve than the primal problem and using the Karush-KuhnTucker conditions we get 
a much simpler problem.

\section{The Karush-KuhnTucker Conditions}